\chapter{Patterns: providers and contexts}
\label{ch:providers}

\begin{aside}
Some state is needed by many widgets but only one source.
Theme, locale, current user, authentication token. Providers
make it ambient.
\end{aside}

\section{The provider}

\begin{lstlisting}
class ThemeProvider {
    Signal<Theme> theme;
public:
    Element provide(Element child) {
        return with_context<ThemeProvider>(*this, child);
    }
};
\end{lstlisting}

A wrapper that injects itself into the context.

\section{The consumer}

\begin{lstlisting}
auto styled = [&] {
    auto& theme = ctx.consume<ThemeProvider>();
    return text("Hello") | fg(theme.colour);
};
\end{lstlisting}

The consumer looks up the provider in the context stack.

\section{Multiple providers}

Multiple providers stack at the root:

\begin{lstlisting}
auto root = theme_provider.provide(
           locale_provider.provide(
           user_provider.provide(
              main_app()
           )));
\end{lstlisting}

\section{Scoped overrides}

A subtree can override a provider by providing a different
one mid-tree. The inner provider shadows the outer.

\section{Recap}

\begin{recap}
\begin{itemize}[leftmargin=*]
  \item Providers inject ambient state.
  \item Consumers look up via type.
  \item Subtrees can override.
\end{itemize}
\end{recap}

\section*{Workshop}

\begin{enumerate}[leftmargin=*]
  \item Build a theme provider.
  \item Consume from a deep widget.
  \item Override the theme in a subtree.
\end{enumerate}
